% Generated draft: requires mathematical and source review.
\documentclass[11pt]{article}
\usepackage[margin=25mm]{geometry}
\usepackage{fontspec}
\setmainfont{Latin Modern Roman}
\usepackage{amsmath,amssymb,mathtools,unicode-math}
\setmathfont{Latin Modern Math}
\usepackage{enumitem,needspace,xurl,hyperref}
\setlength{\parindent}{0pt}
\setlength{\parskip}{5pt}
\setlength{\emergencystretch}{3em}
\newcommand{\N}{\mathbb{N}}
\newcommand{\Z}{\mathbb{Z}}
\newcommand{\Q}{\mathbb{Q}}
\newcommand{\R}{\mathbb{R}}
\newcommand{\C}{\mathbb{C}}
\newcommand{\F}{\mathbb{F}}
\newcommand{\A}{\mathbb{A}}
\newcommand{\PP}{\mathbb{P}}
\newcommand{\E}{\mathbb{E}}
\newcommand{\eps}{\varepsilon}
\newcommand{\dd}{\,\mathrm{d}}
\begin{document}
\section*{UnsolvedMath ID: 20000396}
\textbf{Source problem code:} catalogue entry
\subsection{Definitions and mathematical statement}
Fix a prime power \(q\). For a smooth projective geometrically connected curve \(C/\F_q\) of genus \(g\), its zeta function is
\[
 Z(C,t)=\exp\left(\sum_{n\ge1}\#C(\F_{q^n})\frac{t^n}{n}\right)
       =\frac{P_C(t)}{(1-t)(1-qt)}.
\]
Define \(z_q(g)\) to be the number of distinct such rational functions as \(C\) ranges over genus-\(g\) curves over the fixed field \(\F_q\); curves with the same zeta function count once. Smoothness, projectivity and geometric connectedness are the customary curve convention underlying this formulation and the supplied research report.

The legible part of Question 18 asks how many different zeta functions occur for fixed \(q\) and genus \(g\). A relation symbol between \(g\) and \(q\) is corrupted in the canonical extraction and absent from the displayed statement. Thus a possible intended relative-size or asymptotic restriction on \(g\) is unresolved; the per-genus counting function above does not silently reconstruct that symbol.

The supplied partial-result report states, for every fixed \(q\) and integer \(g\ge0\),
\[
 z_q(g)\le\prod_{n=1}^{g}\bigl(\lfloor4gq^{n/2}\rfloor+1\bigr)
 \le(4g+1)^gq^{g(g+1)/4},
\]
with the empty product equal to one. For \(p=\operatorname{char}\F_q\), suppose \(D=2g/(p-1)\) is an integer. The reported Artin--Schreier construction gives
\[
 z_q(g)\ge\lfloor D/2\rfloor+1-\epsilon_{p,D},\qquad
 \epsilon_{p,D}=\begin{cases}1,&p\text{ odd and }p\mid D+1,\\0,&\text{otherwise}.\end{cases}
\]
The candidate contribution states more specifically that reduced Artin--Schreier curves over this same \(\F_q\) realize \(p\)-ranks \(r(p-1)\) for every integer \(1\le r\le\lfloor D/2\rfloor\), and also rank zero unless \(p\mid D+1\). Here the \(p\)-rank of \(C\) is the integer \(f\) with \(\#\operatorname{Jac}(C)[p](\overline{\F}_q)=p^f\). The report mentions an order-two-pole congruence repair but does not provide equations or construction details. In characteristic two, \(D=2g\), so its stated specialization is \(z_q(g)\ge g+1\) for every \(g\). These are supplied bounds and candidate constructions requiring mathematical review, not an exact determination of \(z_q(g)\) or a resolution of an unspecified asymptotic question.
\subsection{Short English statement}
Count distinct zeta functions of genus-\(g\) curves over a fixed finite field. The supplied report gives an upper bound and characteristic-dependent Artin--Schreier lower bounds, including \(g+1\) in characteristic two. The original relation between \(g\) and \(q\) is corrupted and remains unresolved.
\subsection{Sources}
\begin{itemize}
\item Local UnsolvedMath registry, permanent ID \(20000396\), source code \texttt{catalogue entry}; AIM algebraic number theory corpus, \texttt{aim-algebraic-number-theory-notes.json}, notes[47]. Attribution: Ulam AI / UnsolvedMath Contributors, CC BY 4.0.
\item AIM workshop \emph{Arithmetic statistics over finite fields and function fields}, Question 18: \url{https://aimath.org/pastworkshops/arithstatffieldproblems.pdf}. Original remarks and literature fields are empty; no separate authoritative citation for the reported bounds is supplied.
\item The source labels the statement \texttt{reconstructed\_unverified}. Its verification text identifies an extraction error but ends after an introductory colon; the original string contains an unrendered control character between \(g\) and \(q\).
\item Registry dated August 14, 2026: \texttt{partially\_solved}, \texttt{PARTIAL-PROGRESS}; research attempt 1 records problem status open, partial result, rigor proved, novelty confidence low. The research bounds and candidate construction are machine-generated aids, not peer reviewed, and require independent verification before citation.
\end{itemize}
\end{document}
